% ==============================================================
\section{Introduction}\label{sec:intro}
% ==============================================================

Frequent pattern mining~\cite{agrawal1994fast,han2000mining,fournier2017survey}
is a foundational technique for discovering repeatedly co-occurring itemsets
in transaction data, with broad applications in retail analytics, web mining, and bioinformatics.
Classical methods such as Apriori~\cite{agrawal1994fast},
FP-Growth~\cite{han2000mining}, and ECLAT~\cite{zaki2000scalable}
efficiently enumerate globally frequent patterns,
while sliding-window methods~\cite{dsaa2025} detect locally dense patterns
within specific time intervals.

However, these methods only identify \emph{which} patterns are \emph{when} frequent;
they cannot explain \emph{why} pattern occurrence frequency fluctuates.
In practice, external events such as promotional campaigns, holidays,
regulatory changes, and system updates often drive the temporal variation
of pattern support.
Identifying ``which event caused which pattern to densify, and when''
would enable quantifying campaign effects, detecting side effects,
and diagnosing system anomalies---insights directly relevant to decision-making.

In this work, we formalize this challenge as \textbf{statistical attribution}.
``Attribution'' here refers to identifying the statistical significance
of \textbf{temporal associations} between support change points and external events,
not causal inference.
The output is positioned as a \textbf{screening tool} for further analysis
(see Definition~\ref{def:attribution} and Remark~\ref{rem:attribution_scope}).

Despite its practical importance, no statistical framework exists,
to our knowledge, that links support fluctuations of frequent patterns
to external events.
Temporal pattern mining~\cite{mahanta2005finding,tanbeer2009discovering,kiran2017discovering,fournier2021mining}
discovers patterns with periodic or local frequency but does not model
the association between pattern variation and exogenous events.
Change point detection methods~\cite{page1954continuous,aminikhanghahi2017survey}
identify structural changes in individual time series but lack the capacity
to handle the multiple testing burden of thousands of pattern support series simultaneously
or to link detected changes to events.
The econometric event study method~\cite{mackinlay1997event}
evaluates the impact of corporate events on stock prices but targets a single outcome variable
and cannot be directly applied to thousands of pattern time series simultaneously.

\subsection{Contributions}

The main contributions of this work are a formal definition of the
\textbf{event attribution problem} and the first statistical solution framework.
(1)~We formulate the event attribution problem as a multiple testing problem
with FDR control and identify three challenges: scale, secondary patterns,
and temporal confounding (\S\ref{sec:problem}).
(2)~We construct a five-step pipeline comprising change point extraction,
attribution scoring, circular-shift permutation testing, global BH correction,
and Union-Find deduplication (\S\ref{sec:method}).
(3)~We identify the \emph{secondary pattern problem}---where all patterns containing item $i$
fluctuate jointly when $i$'s frequency changes---and resolve it
via deduplication based on a $\lceil l/2 \rceil$ majority-share criterion (\S\ref{sec:dedup}).
(4)~We provide evaluation across 8 synthetic conditions, ablation analysis,
comparison against 4 related methods, and exploratory application
to real-world data (Dunnhumby) (\S\ref{sec:experiments}).
